
% Pakete zur korrekten Darstellung von Umlauten etc.
\usepackage[utf8]{inputenc}
\usepackage[T1]{fontenc}
\usepackage[ngerman]{babel}
% Paket mit einer schönen Schriftart
\usepackage{palatino}
\usepackage{blindtext}

% Weiter Pakete können hinzugefügt werden
\usepackage{
  booktabs,
  graphicx,
  xcolor,
  ragged2e,
  amsmath,
  amssymb
}

%  Für klickbare Inhaltsverzeichnise, Verweise etc.
\usepackage[
  bookmarks=true,           % Lesezeichen erzeugen
  bookmarksopen=true,
	colorlinks=true,
	linkcolor=black, 
	anchorcolor=black, 
	citecolor=black, 
	filecolor=black, 
	menucolor=black, 
	pagecolor=black,
	urlcolor=black
]{hyperref}

%  Abkürzungsverzeichnis
\usepackage[
     footnote,
     nohyperlinks
]{acronym}

% Neue Makros
\newcommand{\infoprojekt}{Informatikprojekt Klasse 10}

% Für Internetlinks
\usepackage{url}

% Für Quelltexte etc.
\usepackage{listings}
\renewcommand{\lstlistingname}{Programmcode} % adjusts Listings caption and reference lst
\renewcommand{\lstlistlistingname}{Programmcodeverzeichnis}

% Umlaute in Code
% https://tex.stackexchange.com/a/39645
\lstset{literate=%
	{Ö}{{\"O}}1
	{Ä}{{\"A}}1
	{Ü}{{\"U}}1
	{ß}{{\ss}}1
	{ü}{{\"u}}1
	{ä}{{\"a}}1
	{ö}{{\"o}}1
	, tabsize=4
}

% Seiten- und Layouteinstellung
\usepackage[left=3.3cm,%
            top=3cm,%
            bottom=3cm,%
            textwidth=15.5cm,%
            marginparwidth=0cm]{geometry}

\usepackage{setspace}
\onehalfspacing					%%  1,5 zeilig

 
% Für Bildunterschriften etc.
\usepackage{array}
\usepackage[font=small, labelfont=scriptsize,sf,textfont=it]{caption}      

% [H] placement
\usepackage{float}

% Kopf- und Fußzeilen
\usepackage{scrlayer-scrpage}
\clearscrheadings
\clearscrplain
\pagestyle{scrheadings}
\rohead{\rightmark}
\rofoot[\pagemark]{\pagemark}
\setheadsepline{.4pt} % Linie unter dem Head

% Fussnoten bündig
\deffootnote{1.2em}{0em}{%
            \makebox[1.2em][l]{\bfseries\thefootnotemark}}

\usepackage{listings}
\usepackage{xcolor}

\lstset{
    language=Java,
    basicstyle=\scriptsize,
    keywordstyle=\color{blue},
    commentstyle=\color{gray},
    stringstyle=\color{red},
    breaklines=true,
    frame=single
}